% Ad Familiares 13.9 (ad-familiares-13-09)
% Source: https://raw.githubusercontent.com/PerseusDL/canonical-latinLit/master/data/phi0474/phi056/phi0474.phi056.perseus-lat1.xml
% Fetched by scripts/fetch_latin.py

% Dateline (Perseus): Scr., ut videtur, Tarsi vel ex. a. 703 (5 ) vel in. a. 704 (50) .

\ciceroLetterOpener{CICERO CRASSIPEDI S.}

\ciceroSection{1} quamquam tibi praesens commendavi, ut potui diligentissime, socios Bithyniae teque cum mea commendatione tum etiam tua sponte intellexi cupere ei societati quibuscumque rebus posses commodare, tamen cum ii, quorum res agitur, magni sua interesse arbitrarentur me etiam per litteras declarare tibi qua essem erga ipsos voluntate, non dubitavi haec ad te scribere.

\ciceroSection{2} volo enim te existimare me, cum universo ordini publicanorum semper libentissime tribuerim idque magnis eius ordinis erga me meritis facere debuerim, tum in primis amicum esse huic Bithynicae societati, quae societas ordine ipso et hominum genere pars est maxima civitatis (constat enim ex ceteris societatibus), et casu permulti sunt in ea societate valde mihi familiares in primisque is, cuius praecipuum officium agitur hoc tempore, P. Rupilius P. f. Men., qui est magister in ea societate.

\ciceroSection{3} quae cum ita sint, in maiorem modum a te peto Cn. a Pupium, qui est in operis eius societatis, omnibus tuis officiis atque omni liberalitate tueare curesque ut eius operae, quod tibi facile factu est, quam gratissimae sint sociis, remque et utilitatem sociorum (cuius rei quantam potestatem quaestor habeat non sum ignarus) per te quam maxime defensam et auctam velis. id cum mihi gratissimum feceris, tum illud tibi expertus promitto et spondeo, te socios Bithyniae, si iis commodaris, memores esse et gratos cogniturum.
